\question[5]
Von einer Funktion ist nur der Kopf mit dem Docstring bekannt:
\begin{lstlisting}
def kuerze(text, n, ende='...'):
    '''
    text: String
    n: ganze Zahl >= 0
    ende: String, der an einen gekürzten Text angehängt wird
    returns: text, falls er höchstens n Zeichen hat,
             sonst die ersten n Zeichen von text und danach ende
    '''
\end{lstlisting}
\begin{parts}
\part Was ergeben \texttt{kuerze('Informatik', 4)}, \texttt{kuerze('Hallo', 10)} und \texttt{kuerze('Hallo', 2, '!')}?
\part Implementiere die Funktion.
\end{parts}
\begin{solutionbox}{3.25cm}
a) \texttt{'Info...'}, \texttt{'Hallo'}, \texttt{'Ha!'}
\begin{lstlisting}
def kuerze(text, n, ende='...'):
    if len(text) <= n:
        return text
    return text[:n] + ende
\end{lstlisting}
\end{solutionbox}
